% !TEX root = ../thesis.tex

\chapter{Implementácia}
\label{ch:implementation}

Táto kapitola opisuje implementačnú realizáciu tréningovej fázy doménovej adaptácie modelov. Ťažisko je na konkrétnej konfigurácii tréningu a modelovo špecifických technických úpravách, ktoré boli nevyhnutné pre stabilné doučenie a následnú preferenčnú kalibráciu.

\section{Konfigurácia tréningu}

\subsection{Spoločné parametre pipeline}

Pre všetky modely je implementovaný spoločný SFT postup:
\begin{itemize}
    \item tréningový framework: \texttt{TRL SFTTrainer},
    \item dátový formát: JSONL \texttt{prompt/completion},
    \item dĺžka sekvencie: \texttt{max\_seq\_length = 1024},
    \item kvantizácia: \texttt{load\_in\_4bit=True}, typ \texttt{nf4},
    \item optimalizátor: \texttt{paged\_adamw\_8bit},
    \item zapnutý \texttt{gradient\_checkpointing},
    \item uloženie výsledku: adaptér do adresára \texttt{final\_adapter} po ukončení tréningu.
\end{itemize}

Súčasťou pipeline je aj jednotný systémový prompt, ktorý definuje požadovanú štruktúru a kvalitu odpovede (stručnosť, faktickosť, práca s číselnými údajmi, risk a confidence).

\subsection{Hyperparametre podľa modelov}

\textbf{Phi-3 Mini (SFT)}
\begin{itemize}
    \item počet epoch: \texttt{10},
    \item \texttt{per\_device\_train\_batch\_size}: \texttt{1};
    \item \texttt{gradient\_accumulation\_steps}: \texttt{16};
    \item learning rate: \texttt{2e-4};
    \item \texttt{test\_size}: \texttt{0.01};
    \item compute dtype: \texttt{bfloat16};
    \item LoRA: \texttt{r=64}, \texttt{alpha=128}, \texttt{target\_modules=[qkv\_proj, o\_proj, gate\_up\_proj, down\_proj]};
    \item \texttt{warmup\_ratio=0.03}, \texttt{weight\_decay=0.01}.
\end{itemize}

\textbf{Gemma 3 4B (SFT)}
\begin{itemize}
    \item počet epoch: \texttt{3},
    \item \texttt{per\_device\_train\_batch\_size}: \texttt{1};
    \item \texttt{gradient\_accumulation\_steps}: \texttt{8};
    \item learning rate: \texttt{2e-4};
    \item \texttt{test\_size}: \texttt{0.05};
    \item compute dtype: \texttt{bfloat16};
    \item \texttt{bnb\_4bit\_use\_double\_quant=True};
    \item LoRA: \texttt{r=16}, \texttt{alpha=32}, \texttt{target\_modules=[q\_proj, k\_proj, v\_proj, o\_proj, gate\_proj, up\_proj, down\_proj]};
    \item \texttt{max\_grad\_norm=0.3}, \texttt{warmup\_steps=10}, \texttt{lr\_scheduler\_type=constant}.
\end{itemize}

\textbf{Qwen3 4B (SFT)}
\begin{itemize}
    \item počet epoch: \texttt{3},
    \item \texttt{per\_device\_train\_batch\_size}: \texttt{1};
    \item \texttt{gradient\_accumulation\_steps}: \texttt{16};
    \item learning rate: \texttt{2e-4};
    \item \texttt{test\_size}: \texttt{0.01};
    \item compute dtype: \texttt{bfloat16};
    \item LoRA: \texttt{r=16}, \texttt{alpha=32}, \texttt{target\_modules=[q\_proj, k\_proj, v\_proj, o\_proj, gate\_proj, up\_proj, down\_proj]};
    \item \texttt{warmup\_ratio=0.03}, \texttt{weight\_decay=0.01}.
\end{itemize}

\subsection{Modelovo špecifické implementačné detaily}

Okrem hyperparametrov obsahuje pipeline aj niekoľko technických úprav potrebných pre stabilný tréning:
\begin{itemize}
    \item \textbf{Gemma 3 4B:} pred tréningom je aplikovaný patch metódy \texttt{forward}, ktorý pri text-only batchoch dopĺňa \texttt{token\_type\_ids}. Bez tejto úpravy môže multimodálny model zlyhať na čisto textovom vstupe.
    \item \textbf{Gemma 3 4B:} v \texttt{SFTTrainer} sa explicitne používa \texttt{processing\_class=tokenizer}, aby sa obišla automatická voľba multimodálneho procesora.
    \item \textbf{Qwen3 4B:} počas inferencie je \texttt{enable\_thinking=False}, aby výstup neobsahoval interné \enquote{thinking} bloky a bol konzistentný s požadovaným formátom odpovede.
\end{itemize}

\subsection{Dodatočná fáza DPO}

V projekte sa po SFT používa aj dodatočná preference-optimalizácia (DPO) pre všetky tri modely (Phi-3, Gemma 3 a Qwen3):
\begin{itemize}
    \item inicializácia policy modelu zo SFT adaptéra,
    \item referenčný model je odvodený z rovnakého SFT checkpointu,
    \item tréning prebieha na preference datasete vo formáte \texttt{prompt/chosen/rejected},
    \item typická konfigurácia v projekte: \texttt{beta=0.1}, \texttt{epochs=1}, \texttt{learning\_rate=5e-5}, \texttt{gradient\_accumulation\_steps=8}, \texttt{max\_length=1024}, \texttt{max\_prompt\_length=512}, \texttt{eval\_ratio=0.02}.
\end{itemize}

Táto fáza nenahrádza SFT, ale slúži ako mechanizmus dodatočnej kalibrácie preferencií modelu medzi \enquote{lepšou} a \enquote{horšou} odpoveďou na rovnaký vstup.

Implementačne je DPO realizované nasledovne: pre referenčný aj policy model sa najprv načíta kvantizovaná báza modelu, následne sa zlúči SFT adaptér (\texttt{merge\_and\_unload}) a na policy vetvu sa opätovne aplikuje LoRA. Tento postup zodpovedá praxi, kde DPO jemne dolaďuje preferencie už doménovo adaptovaného modelu namiesto tréningu od nuly.
